\begin{helper}[The two balanced blocks at order nine]
Let $(T,e,\varphi,N)$ be a balanced normal globally maximum triple in a
finite loopless multigraph $G$, with $t(G)=9$ and $h(G)=5$.  Then there is
a balanced normal globally maximum triple $(S,e',\psi,M)$ with
$S.\mathsf{length}=8$ and $M.\mathsf{height}=5$, and there are two
distinct colors $\gamma_0,\gamma_1$, such that, writing the positions of
$S$ as
\[
 (y_0,a_0,y_1,a_1,y_2,a_2,y_3,a_3,y_4,
   f_{0,0},u_{0,0},f_{0,1},u_{0,1},
   f_{1,0},u_{1,0},f_{1,1},u_{1,1}),
\]
the following hold for $i,j\in\{0,1\}$:
\begin{enumerate}
\item $f_{i,j}$ has exact endpoint pair $\{y_q,u_{i,j}\}$ for some
      $q\in\{0,1,2,3,4\}$;
\item $\psi(f_{i,j})=\gamma_i$;
\item $\gamma_i$ is missing at some $y_q$ with $0\leq q<5$;
\item the colored base identities $a_1,a_2,a_3$ have colors different
      from both $\gamma_0$ and $\gamma_1$.
\end{enumerate}
All displayed edges are individual edge identities; parallel identities
are not identified.
\end{helper}
\begin{proof}
The ambient object is the finite edge-identity model of
\noderef{LooplessMultigraph}; incidences below are always those of
\noderef{Incident}.  The hypotheses say that $N$ is a
\noderef{NormalGlobalMaximumTriple} and that it satisfies
\noderef{BalancedGlobalMaximumTriple}.  Since its height is five, the
eight displayed edge positions split after the five-vertex normal trunk
into two equal-colored pairs.  Call their colors initially $\gamma_0$
and $\gamma_1$.  The \noderef{TashkinovTree} field supplies a missing
witness for $\gamma_0$ among the five earlier trunk vertices and a
missing witness for $\gamma_1$ among the first seven vertices.  The
global field is \noderef{GlobalMaximumTriple}; it includes a
\noderef{KTriple}, hence the unique-uncolored-edge condition
\noderef{IsOneEdgeUncolored}.  Its order equality, $t(G)=9$, and vertex
injectivity in \noderef{TreeSequence} show that
\noderef{TreeSequenceVertexSet} consists of exactly the nine distinct
displayed vertices.

Choose the first trunk vertex missing $\gamma_0$.  If it is not the
first displayed vertex, apply \noderef{ScheideTrunkRotationPreserves}.
The exact formulas in \noderef{TrunkRotation} move that witness to the
first position, preserve every post-trunk vertex and edge identity, and
alter only the two normal trunk colors on the closing cycle.  To check
that the missing color really survives, use the elementary clause of
\noderef{ScheideGlobalMaximumOrder}, expanded as
\noderef{ColoredVertexSetElementary}: distinct displayed vertices have
disjoint missing-color sets.  If $\gamma_0$ is the first normal color,
its witness was already the first vertex; if it is the second normal
color, its witness is the second vertex and the rotation changes only
the adjacent uncolored and first-normal-color identities; otherwise no
identity whose color is $\gamma_0$ changes.  Thus after rotation
$\gamma_0$ is missing at $y_0$, while both tail pairs and balancedness
are preserved.

From now on write $\alpha,\beta$ for the two current normal trunk colors.
The trunk-edge endpoint clause of
\noderef{NormalGlobalMaximumTriple}, together with its alternating closing
path and properness, gives the exact identities
\[
 a_1:y_1y_2\text{ of color }\alpha,\quad
 a_2:y_2y_3\text{ of color }\beta,\quad
 a_3:y_3y_4\text{ of color }\alpha,
\]
and a closing identity $q:y_4y_0$ of color $\beta$.  The same structure
says that $\alpha$ is missing at $y_0$ and $\beta$ at $y_1$.

The first identity of the first tail pair crosses from the trunk to its
new vertex by the attachment clause of \noderef{TreeSequence}.  Apply
\noderef{ScheideHeightFiveGammaEdges}.  It produces three distinct exact
roles: a $\gamma_0$-identity from $y_1$ to an outside-trunk vertex $v_1$,
a $\gamma_0$-identity from $y_4$ to an outside-trunk vertex $v_2$, and a
$\gamma_0$-identity with endpoint pair $\{y_2,y_3\}$.  Both identities
of the first balanced pair attach to the trunk.  Indeed, the second
cannot attach to the first new vertex, since the two distinct identities
would then be adjacent in the sense of \noderef{AdjacentEdges}, have the
same color, and contradict \noderef{IsProperEdgeColoring}.  Neither can
attach at $y_0$, where the color is missing, nor at $y_2$ or $y_3$, where
the internal identity already gives the unique proper $\gamma_0$
incidence.  Hence their attachment vertices are $y_1$ and $y_4$, one
each.  Properness identifies them with the two crossing identities just
obtained, preserving their individual edge identities.  Interchanging
the two complete edge--new-vertex entries if necessary preserves every
\noderef{TreeSequence} and \noderef{TashkinovTree} condition.  We may
therefore label them
\[
 f_{0,0}=y_1v_1,\qquad f_{0,1}=y_4v_2.
\]

Now apply \noderef{ScheideHeightFiveOutsideLinks}.  It gives two distinct
parallel identities on the exact endpoint pair $\{v_1,v_2\}$, one of
each normal trunk color.  Together with the alternating normal closing
path and the three $\gamma_0$-identities above, these identities form
three matchings on the first seven vertices.  Each matching leaves only
its displayed missing vertex unmatched.  If $\gamma_1$ were either
normal trunk color or $\gamma_0$, neither of its two distinct tree
identities could attach to an earlier vertex: every nonmissing earlier
vertex already has a unique incident identity of that color, while the
missing vertex has none.  This contradicts the attachment field of
\noderef{TreeSequence}.  Thus $\gamma_1$ differs from both trunk colors
and from $\gamma_0$.

The order lemma \noderef{ScheideGlobalMaximumOrder} says that the displayed
vertex set is elementary and closed in the precise senses of
\noderef{ColoredVertexSetElementary} and
\noderef{ColoredVertexSetClosed}.  We use this to obtain three separately
named identities between the last two vertices.  Fix one of the pairwise
distinct colors $c\in\{\gamma_0,\alpha,\beta\}$.  Its displayed matching
on the first seven vertices has exactly one missing vertex.  Elementarity
therefore makes $c$ present at both last vertices, and
\noderef{PresentColors} supplies a $c$-colored identity $g_c$ incident
with the first one.  By \noderef{Incident}, write its endpoint pair as
$\{u_{1,0},w_c\}$.  The distinct-end law of
\noderef{LooplessMultigraph} gives $w_c\ne u_{1,0}$.  If $w_c$ were
outside the tree, this endpoint equality would put $g_c$ in the
\noderef{BoundaryEdges} of the displayed vertex set.  Since $c$ is
missing at its unmatched earlier vertex, however,
\noderef{ColoredVertexSetClosed} forbids a $c$-colored boundary identity.
Thus $w_c$ lies in the tree.

By the equality defining \noderef{TreeSequenceVertexSet}, the tree
vertices are exactly the nine displayed vertices.  At the unique earlier
vertex missing $c$, \noderef{MissingColors} forbids a $c$-incidence.  At
each of the other first seven vertices, properness makes its displayed
$c$-matching identity the unique $c$-identity incident there, and that
identity has both endpoints among the first seven vertices.  Hence $w_c$
cannot be any of those seven vertices, and looplessness excludes the
first last vertex itself.  It must be the second last vertex.  Applying
the argument to $c=\gamma_0,\alpha,\beta$ gives exact identities
$g_{\gamma_0},g_\alpha,g_\beta$ on the common endpoint pair
$\{u_{1,0},u_{1,1}\}$.  Their colors make them three distinct parallel
edge identities.

The witness for $\gamma_1$ lies either on the trunk or among
$v_1,v_2$.  In the latter case let $b_v$ be the old $\beta$-colored
outside-link identity on $\{v_1,v_2\}$ and replace the displayed order by
\[
 (y_0,e,y_1,f_{0,0},v_1,b_v,v_2,f_{0,1},y_4,
   a_1,y_2,a_3,y_3,f_{1,0},u_{1,0},f_{1,1},u_{1,1}),
\].
The endpoint equalities show successively that every new vertex attaches
to an earlier vertex.  The eight identities are distinct by old tree-edge
injectivity, their exact endpoint pairs, and, where an endpoint pair is
shared, their different colors.  The missing witnesses are explicit: the
old $\gamma_0$-witness supports $f_{0,0},f_{0,1}$, the old second
trunk-color witness supports $b_v$, the old first-trunk-color witness
supports $a_1,a_3$, and the assumed old $\gamma_1$-witness in
$\{v_1,v_2\}$ precedes the unchanged last pair.  Thus all clauses of
\noderef{TashkinovTree} hold.  The coloring and unique uncolored identity
are unchanged and the vertex set is the same nine-element set, so the
\noderef{KTriple} and \noderef{GlobalMaximumTriple} clauses hold.  The
first five vertices lie on the simple alternating path
$y_1,f_{0,0},v_1,b_v,v_2,f_{0,1},y_4,q,y_0$, where $q$ is the old normal
closing identity.  This proves \noderef{NormalGlobalMaximumTriple} of
height five, and the remaining pairs $a_1,a_3$ and $f_{1,0},f_{1,1}$
prove \noderef{BalancedGlobalMaximumTriple}.

Now relabel positions and colors, keeping the following old-to-new table
explicit.  The new trunk colors are the old $\gamma_0$ and the old second
trunk color $\beta$; call them the new $\alpha$ and $\beta$.  The new
first-block color is the old first trunk color $\alpha$; call it the new
$\gamma_0$.  The new second-block color is the old $\gamma_1$, still
called $\gamma_1$.  The new first-block outside vertices are the old
$y_2,y_3$.  Their new-$\alpha$ outside link is the old internal
$\gamma_0$-identity on $\{y_2,y_3\}$, while their new-$\beta$ outside
link is the old normal identity $a_2$ on that same endpoint pair.  The
last-pair identities also transport exactly: old $g_{\gamma_0}$ is the
new $g_\alpha$, old $g_\alpha$ is the new $g_{\gamma_0}$, and old
$g_\beta$ is the new $g_\beta$; all retain the exact endpoint pair of the
unchanged last two vertices.  Thus every outside-link and last-pair
identity used below has been checked under the relabeling, and the old
$\gamma_1$-witness is now a trunk witness.  If it was already on the
trunk, no replacement or relabeling is made.  In either branch we
henceforth use the current names $\alpha,\beta,\gamma_0,\gamma_1$.

All the hypotheses of
\noderef{P2968_T2_ScheideOrderNineBlockCompletion} now hold for the
current presentation.  Its first-block identities are exactly the
crossings just named, in the displayed orientation or its reverse; its
second-normal-color link is the current identity on the first outside
pair.  The three last-pair identities, the two trunk missing witnesses,
the color exclusions, and the second-pair coloring were verified above,
including the explicit old-to-new table in the replacement branch.
The graph height is still $h(G)=5$.  Apply the helper to this presentation.
It contains both the seven-vertex Kempe-path contradiction and the
height-five localization, and returns the full block-structure conclusion
for the same coloring and tree.  This supplies all four asserted
properties and completes the proof.
\end{proof}
